\section*{Algorithms}

This appendix collects the pseudocode for the algorithms implementing the item screening procedure described in Section~\ref{sec:it_sec}, in the order in which they are used.

\subsection*{Step 2: Partial Gamma Coefficient}
$A$ and $B$ are variables that correspond to distinct items if we check for LD. If we check for DIF, one is an item and one is an exogenous variable. All rejected tests are flagged and we move onto the next step.

\begin{breakablealgorithm}
\caption{Partial $\gamma$ coefficient for testing $A \bot B \mid \mathcal{C}$}
\label{alg:part_gam}

\begin{algorithmic}
\Require Dataset with item responses and exogenous covariates; variables $A$ and $B$; conditioning variables $\mathcal{C}$
\Ensure Partial gamma coefficient $\gamma_{A \bot B \mid \mathcal{C}}$

\State Initialize
\[
C \gets 0,
\qquad
D \gets 0.
\]

\State Stratify the dataset according to the joint levels of the conditioning variable(s) $\mathcal{C}$.
\State Let $n$ denote the number of resulting strata.

\For{$l = 1,\ldots,n$}

    \State Let $\mathcal{D}_l$ denote the observations in stratum $l$.

    \If{$\mathcal{D}_l$ contains fewer than two observations}
        \State Skip stratum $l$.
    \EndIf

    \State Construct the contingency table of $A$ and $B$ within stratum $l$.

    \If{the contingency table has fewer than two rows or fewer than two columns}
        \State Skip stratum $l$.
    \EndIf

    \State Compute the number of concordant pairs in stratum $l$:
    \[
    C_l \gets \#\{\text{concordant pairs in stratum } l\}.
    \]

    \State Compute the number of discordant pairs in stratum $l$:
    \[
    D_l \gets \#\{\text{discordant pairs in stratum } l\}.
    \]

    \State Update
    \[
    C \gets C + C_l,
    \qquad
    D \gets D + D_l.
    \]

\EndFor

\If{$C + D = 0$}
    \State \Return $\mathrm{NA}$
\EndIf

\State Compute
\[
\gamma_{A \bot B \mid \mathcal{C}}
\gets
\frac{C-D}{C+D}.
\]

\State \Return $\gamma_{A \bot B \mid \mathcal{C}}$

\end{algorithmic}
\end{breakablealgorithm}

\subsection*{Step 3(a): Elimination of Spurious Local Dependence}
\begin{breakablealgorithm}
\caption{Elimination of Spurious Local Dependence}
\label{alg:LD_eli_algo}

\begin{algorithmic}[1]
\Require Set of significant local dependence hypotheses from the initial LD screening.
For each item pair $\{Y_i,Y_{i'}\}$, at most two directional hypotheses can be significant:
\[
H_{ii'}^{(i)}:\; Y_i \perp Y_{i'} \mid R_i,
\qquad
H_{ii'}^{(i')}:\; Y_i \perp Y_{i'} \mid R_{i'}.
\]
Let $\gamma_{ii'}^{(i)}$ and $\gamma_{ii'}^{(i')}$ denote the corresponding partial
$\gamma$ coefficients, and let $w_{ii'}^{(i)}$ and $w_{ii'}^{(i')}$ denote the
corresponding numbers of comparable pairs.

\Ensure Reduced set $\mathcal{G}$ of item pairs considered to exhibit genuine local dependence.

\State Initialize the set of active significant directional LD hypotheses:
\[
\mathcal{H}
\gets
\left\{
H_{ii'}^{(k)}
:
H_{ii'}^{(k)} \text{ is significant},
\;
k \in \{i,i'\}
\right\}.
\]

\State Initialize the set of genuine LD relations:
\[
\mathcal{G} \gets \emptyset.
\]

\State For each item pair $\{Y_i,Y_{i'}\}$ represented in the initial LD
screening, compute the pair-level weighted partial gamma
\[
WPG_{ii'}
=
\frac{
w_{ii'}^{(i)}\gamma_{ii'}^{(i)}
+
w_{ii'}^{(i')}\gamma_{ii'}^{(i')}
}{
w_{ii'}^{(i)} + w_{ii'}^{(i')}
}.
\]

\While{$\mathcal{H} \neq \emptyset$}

    \State Identify all unordered item pairs $\{Y_i,Y_{i'}\}$ represented by at least one
    hypothesis in the current active set $\mathcal{H}$.

    \State Using only the currently active significant directional hypothesis, partition
    the represented item pairs into four evidence lists:
    \[
    \mathcal{L}_{++}^{(2)}
    =
    \left\{
    \{Y_i,Y_{i'}\}
    :
    H_{ii'}^{(i)} \in \mathcal{H},\;
    H_{ii'}^{(i')} \in \mathcal{H},\;
    \gamma_{ii'}^{(i)} > 0,\;
    \gamma_{ii'}^{(i')} > 0
    \right\},
    \]
    \[
    \mathcal{L}_{--}^{(2)}
    =
    \left\{
    \{Y_i,Y_{i'}\}
    :
    H_{ii'}^{(i)} \in \mathcal{H},\;
    H_{ii'}^{(i')} \in \mathcal{H},\;
    \gamma_{ii'}^{(i)} < 0,\;
    \gamma_{ii'}^{(i')} < 0
    \right\},
    \]
    \[
    \mathcal{L}_{+}^{(1)}
    =
    \left\{
    \{Y_i,Y_{i'}\}
    :
    \text{exactly one of }
    H_{ii'}^{(i)},H_{ii'}^{(i')}
    \text{ is in } \mathcal{H}
    \text{ and its partial } \gamma \text{ is positive}
    \right\},
    \]
    \[
    \mathcal{L}_{-}^{(1)}
    =
    \left\{
    \{Y_i,Y_{i'}\}
    :
    \text{exactly one of }
    H_{ii'}^{(i)},H_{ii'}^{(i')}
    \text{ is in } \mathcal{H}
    \text{ and its partial } \gamma \text{ is negative}
    \right\}.
    \]

    \State First, consider pairs with two active significant partial correlations:
    \[
    \mathcal{L}^{(2)}
    =
    \mathcal{L}_{++}^{(2)}
    \cup
    \mathcal{L}_{--}^{(2)}.
    \]

    \If{$\mathcal{L}^{(2)} \neq \emptyset$}
        \State Select the pair with the largest absolute weighted partial gamma among
        pairs with two active significant partial correlations:
        \[
        \{Y_a,Y_b\}
        =
        \arg\max_{\{Y_i,Y_{i'}\}\in \mathcal{L}^{(2)}}
        \left|WPG_{ii'}\right|.
        \]
    \Else
        \State Select the pair with the largest absolute weighted partial gamma among
        pairs with only one active significant partial correlation:
        \[
        \{Y_a,Y_b\}
        =
        \arg\max_{\{Y_i,Y_{i'}\}\in
        \mathcal{L}_{+}^{(1)}\cup\mathcal{L}_{-}^{(1)}}
        \left|WPG_{ii'}\right|.
        \]
    \EndIf

    \State Add the selected item pair to the set of genuine LD relations:
    \[
    \mathcal{G}
    \gets
    \mathcal{G} \cup \{\{Y_a,Y_b\}\}.
    \]

    \State Remove from the active hypothesis set, all directional LD hypotheses whose
    rest-score is defined by either of the selected items:
    \[
    \mathcal{H}
    \gets
    \mathcal{H}
    \setminus
    \left\{
    H_{ii'}^{(k)} \in \mathcal{H}
    :
    k \in \{a,b\}
    \right\}.
    \]

    \State Thus, if an item pair previously had two active significant directional
    hypotheses, but only one of them remains after this removal, the pair may be
    reconsidered in a later iteration in either
    $\mathcal{L}_{+}^{(1)}$ or $\mathcal{L}_{-}^{(1)}$.

\EndWhile

\State \Return $\mathcal{G}$

\end{algorithmic}
\end{breakablealgorithm}

\subsection*{Step 3(b): Elimination of Spurious DIF Sources}
\begin{breakablealgorithm}
\caption{Elimination of Spurious DIF Sources for Item $Y_i$}
\label{alg:step3b}

\begin{algorithmic}[1]
\Require Dataset with item responses and exogenous covariates; item $Y_i$; item set $\mathcal{I}$; initial source set $\mathrm{SOURCE}(Y_i)$; and significance level $\alpha$

\Ensure Final set of genuine DIF sources $\widehat{\mathrm{SOURCE}}(Y_i)$.

\State Compute the total score
\[
S = \sum_{Y_k \in \mathcal{I}} Y_k .
\]

\State Initialize
\[
\mathcal{S}_i^{(0)} \gets \mathrm{SOURCE}(Y_i).
\]

\State Set $t \gets 0$.

\Repeat

    \State Let
    \[
    \mathcal{P}_i^{(t)} \gets \mathcal{S}_i^{(t)},
    \]
    denote the source set used during the current pass.

    \State Set
    \[
    X_{\mathrm{drop}} \gets \emptyset,
    \qquad
    p_{\mathrm{drop}} \gets -\infty .
    \]

    \ForAll{$X_j \in \mathcal{P}_i^{(t)}$}

        \State Use Algorithm~\ref{alg:part_gam} to compute the asymptotic partial gamma test of
        \[
        H_0:
        Y_i \perp X_j
        \mid
        S,\,
        \mathcal{P}_i^{(t)} \setminus \{X_j\}.
        \]

        \State Let $p_j^{(t)}$ denote the corresponding $p$-value.

        \If{$p_j^{(t)} > \alpha$ and $p_j^{(t)} > p_{\mathrm{drop}}$}
            \State Set
            \[
            X_{\mathrm{drop}} \gets X_j,
            \qquad
            p_{\mathrm{drop}} \gets p_j^{(t)} .
            \]
        \EndIf

    \EndFor

    \If{$X_{\mathrm{drop}} = \emptyset$}
        \State Stop the elimination procedure.
    \Else
        \State Remove the non-significant source with the highest $p$-value found in the current pass:
        \[
        \mathcal{S}_i^{(t+1)}
        \gets
        \mathcal{S}_i^{(t)}
        \setminus
        \{X_{\mathrm{drop}}\}.
        \]

        \State Set $t \gets t+1$.
    \EndIf

\Until{no source is removed or $\mathcal{S}_i^{(t)} = \emptyset$}

\State Set
\[
\widehat{\mathrm{SOURCE}}(Y_i)
\gets
\mathcal{S}_i^{(t)} .
\]

\State \Return $\widehat{\mathrm{SOURCE}}(Y_i)$

\end{algorithmic}
\end{breakablealgorithm}

\subsection*{Step 4: Structural Screening of Covariates}
\begin{breakablealgorithm}
\caption{Step 4 structural screening}
\label{alg:step4}
\begin{algorithmic}[1]
\Require Dataset with item responses and candidate covariates;
significance level \(\alpha\)
\Ensure Final set of covariates which satisfy criterion validity.
\State Compute the total score \(S\) as the row sum of the item responses.
\State Let \(\mathcal{C}\) be the set of candidate covariates.
\State Let \(\mathcal{R} \gets \emptyset\) be the set of removed covariates.

\While{\(\mathcal{C}\) is not empty}
    \For{each covariate \(X_j \in \mathcal{C}\)}
        \State Test
        \[
            H_0: S \perp X_j \mid \mathcal{C} \setminus \{X_j\},
        \]
        using an asymptotic conditional partial gamma test
        \State Store the partial gamma estimate and the corresponding
        \(p\)-value.
    \EndFor

    \If{all decision \(p\)-values are less than or equal to \(\alpha\)}
        \State Stop and retain all covariates in \(\mathcal{C}\).
    \Else
        \State Select the covariate \(X_j\) with the highest decision
        \(p_j > \alpha\).
        \State Remove \(X_j\) from \(\mathcal{C}\).
        \State Add \(X_j\) to \(\mathcal{R}\).
    \EndIf
\EndWhile

\State \Return retained covariates \(\mathcal{C}\) and removed covariates
\(\mathcal{R}\)
\end{algorithmic}
\end{breakablealgorithm}
